% ==============================================================================
% CHƯƠNG 2: HIỂU DỮ LIỆU VÀ EDA (DATA UNDERSTANDING & EDA)
% ==============================================================================

\chapter{Hiểu Dữ Liệu và Khám Phá EDA (Data Understanding)}
\label{chap:data_understanding}

\section{Nguồn Thu Thập Dữ Liệu Nghiên Cứu Khoa Học}

Để xây dựng một kho dữ liệu học thuật quy mô lớn, có tính cập nhật thời sự và bao quát cả các công trình nền tảng lịch sử, hệ thống khai thác dữ liệu từ sáu nguồn học thuật uy tín:

\begin{enumerate}[label=\textbf{\arabic*.}]
  \item \textbf{arXiv RSS Feeds \& OAI-PMH API}: Nguồn dữ liệu chủ lực cung cấp các công trình nghiên cứu mới nhất thuộc các nhóm ngành Khoa học Máy tính (Computer Science) và Thống kê (Statistics). Hệ thống sử dụng giao thức OAI-PMH với cơ chế phân trang tự động \texttt{resumptionToken} để thu thập hàng chục nghìn bài báo kèm siêu dữ liệu chuẩn Dublin Core.
  \item \textbf{arXiv HTML5 / ar5iv Project}: Nguồn văn bản giàu định dạng (rich document source) do dự án ar5iv (hợp tác cùng nhóm LaTeXML tại Viện Tiêu chuẩn và Kỹ thuật Quốc gia Hoa Kỳ - NIST) chuyển đổi trực tiếp từ mã nguồn TeX của tác giả sang định dạng HTML5 ngữ nghĩa. Đây là nguồn dữ liệu duy nhất cho phép bóc tách cấu trúc công thức toán học và các thẻ phân đoạn bài báo với độ chính xác tuyệt đối.
  \item \textbf{OpenAlex \& Semantic Scholar}: Nguồn bổ trợ cung cấp đồ thị trích dẫn học thuật toàn cầu, chỉ số ảnh hưởng của tác giả, và các công trình kinh điển (Foundation Papers như \textit{Attention Is All You Need}, \textit{ResNet}, \textit{BERT}, \textit{Adam Optimizer}) giúp xây dựng mạng lưới trích dẫn xương sống.
  \item \textbf{OpenReview.net}: Thu thập các bài báo hội nghị (ICLR, NeurIPS) cùng với toàn bộ đánh giá của phản biện (peer review critiques, scores, rebuttal) giúp hệ thống RAG có khả năng tra cứu các góc nhìn phản biện học thuật đa chiều.
  \item \textbf{Computer Vision Foundation (CVF) Open Access}: Bổ sung các bài báo hội nghị chuyên ngành Thị giác máy tính đỉnh cao như CVPR 2024.
  \item \textbf{Zenodo Open Repository}: Lưu trữ các tập dữ liệu thực nghiệm và tài liệu bổ trợ khoa học.
\end{enumerate}

\section{Cấu Trúc Bản Ghi và Siêu Dữ Liệu Học Thuật (Silver Schema)}

Dữ liệu sau khi thu thập từ các nguồn thô được chuẩn hóa và tổ chức thành các bảng Parquet dạng cột trong tầng Silver. Mỗi bản ghi tài liệu khoa học tuân thủ nghiêm ngặt cấu trúc lược đồ (schema) gồm các trường thông tin sau:

\begin{itemize}
  \item \texttt{paper\_id} (String): Mã định danh duy nhất của bài báo (ví dụ: \texttt{2401.13216}, \texttt{2310.01407}).
  \item \texttt{title} (String): Tiêu đề bài báo đã được làm sạch khoảng trắng và chuẩn hóa ký tự Unicode.
  \item \texttt{abstract} (String): Tóm tắt nội dung bài báo.
  \item \texttt{authors} (List[String]): Danh sách tên đầy đủ của tất cả các tác giả tham gia công trình.
  \item \texttt{primary\_category} (String): Phân ngành chính theo phân loại của arXiv (ví dụ: \texttt{cs.LG}, \texttt{cs.CV}).
  \item \texttt{categories} (List[String]): Danh sách tất cả các nhãn chuyên ngành mà tác giả đăng ký (thể hiện tính liên ngành).
  \item \texttt{published\_date} (Date): Ngày công bố chính thức hoặc ngày cập nhật phiên bản mới nhất.
  \item \texttt{total\_sections} (Integer): Tổng số mục/chương có cấu trúc được bóc tách từ bài báo.
  \item \texttt{total\_math\_count} (Integer): Tổng số lượng công thức toán học dạng LaTeX ($...$ và $$...$$).
  \item \texttt{total\_words} (Integer): Tổng số lượng từ vựng của toàn văn bài báo.
  \item \texttt{clean\_full\_text} (String): Toàn văn bài báo đã được lọc bỏ các thẻ HTML rác và giữ lại cấu trúc ngữ nghĩa.
  \item \texttt{sections\_json} (JSON String): Cấu trúc phân đoạn chi tiết, bao gồm tiêu đề mục (\texttt{section\_title}), loại mục (\texttt{section\_type}: abstract, introduction, methodology, experiments, conclusion), và nội dung chi tiết.
  \item \texttt{content\_hash} (String): Mã băm SHA-256 bảo đảm toàn vẹn nội dung và phục vụ lọc bỏ trùng lặp.
\end{itemize}

\section{Thống Kê Mô Tả Toàn Cảnh Hồ Dữ Liệu (Descriptive Statistics)}

Dựa trên dữ liệu thực tế được tính toán bởi công cụ phân tích DuckDB OLAP (ghi nhận trong tệp \texttt{eda\_summary.json}), hồ dữ liệu đã xử lý hơn **13,000 bài báo khoa học**, trong đó đợt phân tích chuyên sâu chuẩn hóa ghi nhận **10,000 bài báo** với **9,015 bài báo có toàn văn HTML làm giàu thành công (đạt tỷ lệ làm giàu $90.15\%$)**. 

Toàn bộ kho dữ liệu chứa tổng cộng **47,782,707 từ vựng** (trung bình 4,778 từ/bài) và bảo toàn thành công **2,220,938 công thức toán học LaTeX** (trung bình 222.1 công thức/bài). Khi tính gộp toàn bộ các tầng dữ liệu (bao gồm cả các bài báo lịch sử và hội nghị CVPR/OpenReview), tổng số công thức toán học được bảo toàn lên tới **2,765,395 công thức** trên **60.98 triệu từ vựng**.

Bảng~\ref{tab:eda_categories} trình bày phân bố chi tiết số lượng bài báo, công thức toán và độ dài trung bình theo từng chuyên ngành khoa học.

% ==============================================================================
% BẢNG PHÂN BỐ DỮ LIỆU THEO CHUYÊN NGÀNH & ĐẶC TÍNH TOÁN HỌC
% ==============================================================================

\begin{table}[htbp]
  \centering
  \small
  \caption{Phân bố Số lượng Bài báo, Công thức Toán học và Độ dài Từ vựng theo Chuyên ngành trên Hồ Dữ liệu Học thuật (Nguồn: \texttt{eda\_summary.json})}
  \label{tab:eda_categories}
  \begin{tabularx}{\textwidth}{@{}l c c r r c@{}}
    \toprule
    \textbf{Chuyên ngành (arXiv Category)} & \textbf{Số bài báo} & \textbf{Tỷ lệ (\%)} & \textbf{Tổng công thức toán} & \textbf{Mật độ (ct/bài)} & \textbf{Số từ TB/bài} \\
    \midrule
    \textbf{\texttt{cs.LG}} (Machine Learning) & 2,380 & 23.80\% & 886,955 & 372.7 & 5,099.7 \\
    \textbf{\texttt{cs.CV}} (Computer Vision) & 2,271 & 22.71\% & 308,487 & 135.8 & 4,432.7 \\
    \textbf{\texttt{cs.CL}} (Computation \& Language) & 1,467 & 14.67\% & 149,869 & 102.2 & 4,315.5 \\
    \textbf{\texttt{cs.RO}} (Robotics) & 682 & 6.82\% & 101,023 & 148.1 & 4,636.6 \\
    \textbf{\texttt{cs.AI}} (Artificial Intelligence) & 453 & 4.53\% & 94,302 & 208.2 & 5,383.9 \\
    \textbf{\texttt{eess.IV}} (Image \& Video Processing) & 294 & 2.94\% & 32,010 & 108.9 & 3,559.7 \\
    \textbf{\texttt{cs.IR}} (Information Retrieval) & 253 & 2.53\% & 22,140 & 87.5 & 4,210.3 \\
    \textbf{\texttt{stat.ML}} (Statistical Machine Learning) & 210 & 2.10\% & 142,740 & \textbf{679.7} & 5,420.1 \\
    \textbf{Các chuyên ngành khác} & 1,990 & 19.90\% & 483,412 & 242.9 & 4,120.4 \\
    \midrule
    \textbf{Tổng cộng toàn tập dữ liệu} & \textbf{10,000} & \textbf{100.0\%} & \textbf{2,220,938} & \textbf{222.1} & \textbf{4,778.3} \\
    \bottomrule
  \end{tabularx}
\end{table}


\section{Các Phát Hiện EDA Cốt Lõi (Core Domain Insights)}

Quá trình khám phá dữ liệu thực nghiệm (EDA) kết hợp truy vấn cột DuckDB đã phát hiện những hiểu biết sâu sắc (insights) có ý nghĩa quyết định đối với việc thiết kế hệ thống khai phá và truy xuất RAG:

\subsection{Insight 1: "Đế chế toán học" tập trung ở \texttt{cs.LG} và \texttt{stat.ML}}
Số liệu Bảng~\ref{tab:eda_categories} cho thấy sự bất đối xứng rất lớn về hàm lượng toán học giữa các chuyên ngành:
\begin{itemize}
  \item Xét về \textbf{tổng số lượng công thức}: Chuyên ngành Học máy (\texttt{cs.LG}) chiếm vị trí thống trị tuyệt đối với **886,955 công thức toán** (chiếm tới gần 40\% toàn bộ công thức toán của cả kho dữ liệu). Mặc dù \texttt{cs.LG} và \texttt{cs.CV} có số bài xấp xỉ nhau (~2,300 bài), nhưng lượng công thức toán của \texttt{cs.LG} lại cao gấp **2.87 lần** so với \texttt{cs.CV}.
  \item Xét về \textbf{mật độ toán học trên từng bài báo (Mathematical Density)}: Chuyên ngành Thống kê Học máy (\texttt{stat.ML}) mới thực sự là "ông vua lý thuyết". Trung bình mỗi bài báo \texttt{stat.ML} chứa tới **679.7 công thức toán**, cao gấp 1.8 lần so với \texttt{cs.LG} (372.7 công thức/bài), cao gấp 5.0 lần so với \texttt{cs.CV} (135.8 công thức/bài), và gấp gần 7.0 lần so với chuyên ngành Xử lý ngôn ngữ tự nhiên \texttt{cs.CL} (102.2 công thức/bài).
\end{itemize}

\begin{calloutbox}{Ý nghĩa đối với thiết kế hệ thống}
Phát hiện này định hình chiến lược phân đoạn văn bản (chunking strategy): Không thể áp dụng cùng một kích thước đoạn (chunk size) cố định cho tất cả các chuyên ngành. Đối với các bài báo lý thuyết nặng về toán học (\texttt{stat.ML}, \texttt{cs.LG}), bộ phân đoạn phải có cơ chế nhận biết ranh giới phương trình để tránh cắt đôi một khối chứng minh toán học, trong khi các bài báo ứng dụng (\texttt{cs.CL}, \texttt{cs.CV}) cần ưu tiên độ phủ của ngữ cảnh văn bản tự nhiên.
\end{calloutbox}

\subsection{Insight 2: \texttt{cs.AI} đóng vai trò "Trạm Trung chuyển Liên ngành" (Interdisciplinary Hub)}
Một hiện tượng bất ngờ là chuyên ngành Trí tuệ Nhân tạo (\texttt{cs.AI}) chỉ đứng thứ 5 về số lượng bài báo đơn lẻ (453 bài, chiếm 4.53\% dữ liệu). Tuy nhiên, khi phân tích ma trận đồng xuất bản (\textbf{Co-occurrence Matrix}), nhãn \texttt{cs.AI} lại xuất hiện trong hầu hết các cặp liên ngành lớn nhất:
\begin{itemize}
  \item $\texttt{cs.AI} \times \texttt{cs.LG}$: **1,983 bài** (Cặp liên ngành lớn nhất toàn bộ hồ dữ liệu!).
  \item $\texttt{cs.AI} \times \texttt{cs.CL}$: **1,046 bài** (Giao thoa giữa AI tổng quát và Ngôn ngữ học tính toán).
  \item $\texttt{cs.CV} \times \texttt{cs.LG}$: **984 bài** (Thị giác máy tính kết hợp các kỹ thuật Học sâu).
  \item $\texttt{cs.AI} \times \texttt{cs.CV}$: **845 bài** (Hướng tiếp cận Đa phương thức - Multimodal AI).
\end{itemize}
\textit{Ý nghĩa thực tiễn}: Trong cộng đồng nghiên cứu hiện đại, các tác giả hiếm khi xuất bản một bài báo thuần túy dán nhãn \texttt{cs.AI}. Nhãn \texttt{cs.AI} hoạt động như một "siêu nhãn" (umbrella tag) bao bọc các nghiên cứu ứng dụng thực nghiệm vào Thị giác (\texttt{cs.CV}) hoặc Ngôn ngữ (\texttt{cs.CL}). Đây là cơ sở khoa học để xây dựng luật kết hợp mở rộng truy vấn trong Trụ cột 1.

\subsection{Insight 3: Điểm bùng phát thời gian (Temporal Inflection Point 2023--2024)}
Phân tích chuỗi thời gian ngày công bố ghi nhận:
\begin{itemize}
  \item Giai đoạn từ 2005 đến 2022: Tốc độ gia tăng tài liệu diễn ra từ từ, lũy tiến tuyến tính.
  \item Năm 2023: Tốc độ tăng tốc rõ rệt với gần 1,000 bài báo.
  \item **Tháng 1 và Tháng 2 năm 2024**: Bùng nổ dữ dội với **8,838 bài báo** (chiếm tới 68\% tổng số dữ liệu thu thập).
\end{itemize}
Điểm bùng phát này phản ánh chính xác làn sóng bùng nổ của kỷ nguyên Generative AI, các bài báo chuẩn bị cho các hội nghị đỉnh cao như CVPR 2024, ICML 2024 và ICLR 2024.

\subsection{Insight 4: Phân bố 4 góc phần tư giữa Mật độ Toán học và Độ dài Từ vựng}
Khi vẽ biểu đồ phân tán 2D giữa Số lượng từ vựng ($X \in [0, 12000]$) và Số lượng công thức toán ($Y \in [0, 1200]$), hồ dữ liệu phân rã tự nhiên thành bốn góc phần tư rõ rệt:
\begin{itemize}
  \item \textbf{Góc Q1 (Toán dày, Độ dài ngắn - Theoretical Rigor)}: Tập trung các bài báo \texttt{stat.ML} và \texttt{cs.LG} thuần lý thuyết, chứa nhiều chứng minh hội tụ và bổ đề giải tích.
  \item \textbf{Góc Q2 (Toán dày, Độ dài lớn - Foundational Monographs)}: Các công trình nghiên cứu quy mô lớn, luận án hoặc chuyên khảo nền tảng kết hợp cả lý thuyết sâu lẫn thực nghiệm rộng.
  \item \textbf{Góc Q3 (Toán ít, Độ dài ngắn - Short Communications)}: Các bài báo truyền thông hội nghị ngắn mô tả một phát hiện hoặc ứng dụng hẹp.
  \item \textbf{Góc Q4 (Toán ít, Độ dài lớn - Empirical Systems \& Benchmarks)}: Tập trung chủ yếu các bài báo \texttt{cs.CL} và hệ thống LLMs, chủ yếu trình bày kiến trúc, phân tích bảng kết quả thực nghiệm và ablation study dài hàng chục trang.
\end{itemize}

\section{Vấn đề Chất lượng Dữ liệu và Giải Pháp Xử Lý}

Trong quá trình thu thập và tiền xử lý, hệ thống đối mặt với ba thách thức lớn về chất lượng dữ liệu:
\begin{enumerate}[label=\textbf{\arabic*.}]
  \item \textbf{Mất liên kết LaTeX do lỗi phân tích cú pháp}: Nhiều công thức toán học bị công cụ parse thông thường chuyển thành các ký tự rác. Hệ thống giải quyết bằng cách viết riêng module \texttt{AcademicHTMLParser} dựa trên \texttt{BeautifulSoup} và \texttt{lxml}, ưu tiên bóc tách từ thẻ \texttt{annotation encoding="application/x-tex"} để giữ nguyên vẹn 100\% cú pháp LaTeX.
  \item \textbf{Hiện tượng trùng lặp tài liệu giữa các phiên bản (Version Duplication)}: arXiv thường xuyên lưu trữ nhiều phiên bản (v1, v2, v3) của cùng một bài báo. Hệ thống xử lý bằng thuật toán chuẩn hóa mã bài báo (loại bỏ hậu tố phiên bản \texttt{v\textbackslash d+}) và chỉ giữ lại phiên bản có số lượng mục và từ vựng lớn nhất, kết hợp kiểm tra mã băm nội dung SHA-256.
  \item \textbf{Độ lệch thời gian dồn cục (Collection Spike)}: Hiện tượng tập trung dữ liệu vào tháng 01/2024 do đặc thù đợt cào dữ liệu được nhận diện rõ ràng trong báo cáo nhằm tránh suy diễn sai lệch về xu hướng lịch sử dài hạn.
\end{enumerate}
